
\begin{proof}[引理~\ref{lem:7-4-73} 的证明]
由定义，
\begin{equation}
\begin{aligned}
(\delta Y_a)(\mathbf z)
={}&
\sum_c\pm
\mathfrak m\bigl(
\mathbf z_c^{3;1},
Y_a(\mathbf z_c^{3;2}),
\mathbf z_c^{3;3}
\bigr)\\
&+
\sum_c\pm
Y_a\bigl(
\mathbf z_c^{3;1},
\mathfrak m(\mathbf z_c^{3;2}),
\mathbf z_c^{3;3}
\bigr).
\end{aligned}
\tag{7.4.74}\label{eq:7-4-74}
\end{equation}
将
\[
\mathfrak m
=\mathfrak m_{1,\beta_0}
+\bigl(\mathfrak m-\mathfrak m_{1,\beta_0}\bigr)
\]
分解，
并回顾
$\overline{\mathfrak m}_1
=\mathfrak m_{1,\beta_0}$。
依据这一分解，
用式~\eqref{eq:7-4-74}
拆分引理~\ref{lem:7-4-73}
左端各项。
其中含
$\mathfrak m_{1,\beta_0}$
的项对应于引理右端第二项；
含
$\mathfrak m_{k,\beta}$、
$(k,\beta)\ne(1,\beta_0)$
的项对应于右端第一项。
引理得证。
见图~\ref{fig:7-4-9}。
\end{proof}
\par\medskip\noindent\begin{minipage}{\textwidth}
\centering
\setlength{\unitlength}{0.8mm}
\begin{picture}(155,63)
\linethickness{0.55pt}
\put(25,52){\line(0,-1){10}}
\multiput(25,54)(0,2){4}{\line(0,1){0.8}}
\put(25,42){\line(0,-1){20}}
\qbezier(25,42)(16,36)(7,30)
\qbezier(25,42)(21,36)(17,30)
\qbezier(25,42)(29,36)(33,30)
\qbezier(25,42)(34,36)(43,30)
\qbezier(25,22)(16,16)(7,10)
\qbezier(25,22)(21,16)(17,10)
\qbezier(25,22)(29,16)(33,10)
\qbezier(25,22)(34,16)(43,10)
\multiput(16,8)(-1,-2){3}{\line(0,-1){0.8}}
\multiput(34,8)(1,-2){3}{\line(0,-1){0.8}}
\put(27,45){$\scriptstyle\mathfrak m_{k,\beta}$}
\put(27,32){$\scriptstyle\operatorname{id}$}
\put(27,21){$\scriptstyle Y$}
\put(78,52){\line(0,-1){10}}
\multiput(78,54)(0,2){4}{\line(0,1){0.8}}
\put(78,42){\line(0,-1){20}}
\qbezier(78,42)(69,36)(60,30)
\qbezier(78,42)(74,36)(70,30)
\qbezier(78,42)(82,36)(86,30)
\qbezier(78,42)(87,36)(96,30)
\qbezier(78,22)(69,16)(60,10)
\qbezier(78,22)(74,16)(70,10)
\qbezier(78,22)(82,16)(86,10)
\qbezier(78,22)(87,16)(96,10)
\multiput(69,8)(-1,-2){3}{\line(0,-1){0.8}}
\multiput(87,8)(1,-2){3}{\line(0,-1){0.8}}
\put(80,45){$\scriptstyle Y$}
\put(80,32){$\scriptstyle\operatorname{id}$}
\put(80,21){$\scriptstyle\mathfrak m_{k',\beta'}$}
\put(128,42){\line(0,-1){20}}
\qbezier(128,22)(119,16)(110,10)
\qbezier(128,22)(125,16)(122,10)
\qbezier(128,22)(131,16)(134,10)
\qbezier(128,22)(137,16)(146,10)
\put(112,46){$\scriptstyle
\overline{\mathfrak m}_1Y
\ \pm\ Y\widehat{\overline{\mathfrak m}}_1$}
\put(51,0){\makebox(0,0){第一项}}
\put(128,0){\makebox(0,0){第二项}}
\end{picture}
\par\smallskip
{\renewcommand{\thefigure}{7.4.9}\refstepcounter{figure}
\normalfont\bfseries 图~\thefigure\label{fig:7-4-9}}
\end{minipage}\par\medskip


接下来注意到
\begin{equation}
\begin{aligned}
&\sum_{\Gamma\in G(k,\mathbf Y)}
\sum_e
\pm T^{E(\Gamma)}e^{\mu(\Gamma)/2}
\mathfrak h\bigl(
|\Gamma|,
\partial_e Z_\Gamma^{\mathfrak g,II}
\bigr)(\mathbf x)\\
&\quad=
\sum_c
\sum_{\mathbf Y_{(1)}
\cup\mathbf Y_{(2)}=\mathbf Y}
\pm
\mathfrak g(\mathbf Y_{(1)})
\bigl(
\mathbf x_c^{3;1},
\mathfrak g(\mathbf Y_{(2)})
(\mathbf x_c^{3;2}),
\mathbf x_c^{3;3}
\bigr).
\end{aligned}
\tag{7.4.75}\label{eq:7-4-75}
\end{equation}
见图~\ref{fig:7-4-10}。
由引理~\ref{lem:7-4-64}、
引理~\ref{lem:7-4-71}、
式~\eqref{eq:7-4-72}、
引理~\ref{lem:7-4-73}
和式~\eqref{eq:7-4-75}，
除符号外，
命题~\ref{prop:7-4-56}
已得证。
\par\medskip\noindent
\begin{minipage}{0.49\textwidth}
\centering
\setlength{\unitlength}{0.7mm}
\begin{picture}(88,77)
\linethickness{0.55pt}
\put(43,68){\line(0,-1){8}}
\qbezier(43,60)(27,54)(13,48)
\qbezier(13,48)(10,41)(7,36)
\qbezier(13,48)(17,41)(20,36)
\qbezier(43,60)(36,47)(31,36)
\put(43,60){\line(0,-1){26}}
\qbezier(43,60)(55,53)(58,48)
\qbezier(58,48)(55,41)(53,36)
\qbezier(58,48)(64,41)(69,36)
\qbezier(43,60)(66,54)(79,48)
\qbezier(79,48)(76,41)(73,36)
\qbezier(79,48)(83,41)(85,36)
\qbezier(43,34)(27,26)(15,22)
\qbezier(15,22)(10,14)(6,7)
\put(15,22){\line(0,-1){15}}
\qbezier(15,22)(20,14)(24,7)
\qbezier(43,34)(35,19)(30,7)
\qbezier(43,34)(47,19)(52,7)
\qbezier(43,34)(60,26)(70,22)
\qbezier(70,22)(66,14)(63,7)
\qbezier(70,22)(76,14)(79,7)
\put(63,46){\vector(0,-1){16}}
\put(63,30){\vector(0,1){16}}
\put(57,38){\line(1,0){12}}
\put(47,44){$\scriptstyle\mathfrak G$}
\put(73,62){$\scriptstyle\mathfrak g(\mathbf Y_{(2)})$}
\put(73,25){$\scriptstyle\mathfrak g(\mathbf Y_{(1)})$}
\end{picture}
\par\smallskip
{\renewcommand{\thefigure}{7.4.10}\refstepcounter{figure}
\normalfont\bfseries 图~\thefigure\label{fig:7-4-10}}
\end{minipage}\hfill
\begin{minipage}{0.49\textwidth}
\centering
\setlength{\unitlength}{0.7mm}
\begin{picture}(88,77)
\linethickness{0.55pt}
\put(44,68){\line(0,-1){10}}
\qbezier(44,58)(31,35)(15,7)
\put(44,58){\line(0,-1){28}}
\qbezier(44,30)(38,18)(34,7)
\qbezier(44,30)(49,18)(54,7)
\qbezier(44,58)(60,49)(71,40)
\put(71,40){\line(0,-1){21}}
\qbezier(71,19)(67,13)(65,7)
\qbezier(71,19)(77,13)(81,7)
\put(46,64){$\scriptstyle\Pi_H$}
\put(35,54){$\scriptstyle\mathfrak m_3$}
\put(47,43){$\scriptstyle\mathfrak G^1$}
\put(45,29){$\scriptstyle Y_2$}
\put(58,49){$\scriptstyle\mathfrak G^2$}
\put(73,39){$\scriptstyle Y_1$}
\put(73,31){$\scriptstyle\mathfrak G^3$}
\put(73,18){$\scriptstyle\mathfrak m_2$}
\put(12,2){$\scriptstyle x_1$}
\put(31,2){$\scriptstyle x_2$}
\put(51,2){$\scriptstyle x_3$}
\put(62,2){$\scriptstyle x_4$}
\put(79,2){$\scriptstyle x_5$}
\end{picture}
\par\smallskip
{\renewcommand{\thefigure}{7.4.11}\refstepcounter{figure}
\normalfont\bfseries 图~\thefigure\label{fig:7-4-11}}
\end{minipage}\par\medskip


现在核对符号；
同时也将核对
式~\eqref{eq:7-4-61}
中的符号。
为此，
仿照图~\ref{fig:7-4-5}
附近的讨论，
说明式~\eqref{eq:7-4-68}
的符号约定。
取上方图~\ref{fig:7-4-11}
中的树 $\Gamma$ 为例。
它对应于
\begin{equation}
\begin{aligned}
\mathfrak g(\Gamma)
=\pm\Pi_H\Bigl(
\mathfrak m_3\bigl(
x_1,\,
\mathfrak G^1(Y_2(x_2,x_3)),\,
\mathfrak G^2\bigl(
Y_1\bigl(
\mathfrak G^3(
\mathfrak m_2(x_4,x_5))
\bigr)
\bigr)
\bigr)
\Bigr).
\end{aligned}
\tag{7.4.76}\label{eq:7-4-76}
\end{equation}
此式给三个
$\mathfrak G$ 标上
$1,2,3$，
只是为彼此区分；
它们实际上是同一个算子
$\mathfrak G$。

按照图~\ref{fig:7-4-5}
附近的说明，
把算子移至左侧，
把各 $x_i$ 移至右侧，
得到
\begin{equation}
\Pi_H,\,
\mathfrak m_3,\,
\mathfrak G^1,\,
Y_2,\,
\mathfrak G^2,\,
Y_1,\,
\mathfrak G^3,\,
\mathfrak m_2,\,
x_1,x_2,x_3,x_4,x_5.
\tag{7.4.77}\label{eq:7-4-77}
\end{equation}
这一移动产生符号
$(-1)^{*(\mathfrak h)}$，
其中
\begin{equation}
*(\mathfrak h)
=\deg' x_1\,\deg'Y_2
+\bigl(
\deg' x_1+\deg' x_2+\deg' x_3
\bigr)\deg'Y_1.
\tag{7.4.78}\label{eq:7-4-78}
\end{equation}
注意从式~\eqref{eq:7-4-76}
到式~\eqref{eq:7-4-77}
时，
$x_1,x_2,x_3$
与 $\mathfrak m_2$ 换了次序，
但这一换序未出现在
式~\eqref{eq:7-4-78}
中。
原因在于
$\mathfrak m_k$
总与 $\mathfrak G$ 配对，
且
$\deg\mathfrak m_k
+\deg\mathfrak G$
为偶数；
见式~\eqref{eq:7-4-52}。
又因这里把
$D$ 视作 $C[1]$，
$x_i$ 的次数使用移位次数。
至于 $Y_i$，
与 $x_j$ 换序时，
它总与左邻的
$\mathfrak G$
一同移动；
见式~\eqref{eq:7-4-51}。
因此出现
$\deg'Y_i
=\deg Y_i+\deg\mathfrak G$。

由定义，
\[
\begin{aligned}
&\mathfrak h\bigl(
|\Gamma|,Z_\Gamma^{\mathfrak g}
\bigr)(\mathbf x)\\
&\quad=
(-1)^{*(\mathfrak h)}
\Pi_H\Bigl(
\mathfrak m_3\bigl(
x_1,\,
\mathfrak G^1(Y_2(x_2,x_3)),\,
\mathfrak G^2\bigl(
Y_1\bigl(
\mathfrak G^3(
\mathfrak m_2(x_4,x_5))
\bigr)
\bigr)
\bigr)
\Bigr),
\end{aligned}
\]
其中
$*(\mathfrak h)$
由式~\eqref{eq:7-4-78}
给出。

再对式~\eqref{eq:7-4-77}
重新排列：
这次把
$\mathfrak m$ 移至左侧，
$Y_i$ 移至中间，
$x_i$ 移至右侧，
并按 $i$ 重排 $Y_i$；
同时，
每个 $\mathfrak G$
都随其右邻的
$Y_i$ 或 $\mathfrak m_i$
一起移动。
在本例中，
得到
\begin{equation}
\Pi_H,\,
\mathfrak m_3,\,
\mathfrak G^3,\,
\mathfrak m_2,\,
\mathfrak G^2,\,
Y_1,\,
\mathfrak G^1,\,
Y_2,\,
x_1,x_2,x_3,x_4,x_5.
\tag{7.4.79}\label{eq:7-4-79}
\end{equation}
所得符号为
$(-1)^{*(\Gamma)}$，
其中
$*(\Gamma)
=\deg'Y_1\,\deg'Y_2$，
与式~\eqref{eq:7-4-68}
一致。
与上面同理，
$Y_i$ 与 $\mathfrak m_k$
交换时不产生符号，
$Y_i$ 的次数也因此发生同样的移位。
这一讨论同时说明，
式~\eqref{eq:7-4-61}
中的符号是正确的。

所以，
当从式~\eqref{eq:7-4-76}
转到式~\eqref{eq:7-4-79} 时，
\[
\begin{aligned}
\mathfrak g(\Gamma)(\mathbf x)
={}&
(-1)^{*(\mathfrak h)+*(\Gamma)}
\Pi_H\Bigl(
\mathfrak m_3\bigl(
x_1,\,
\mathfrak G^1(Y_2(x_2,x_3)),\\
&\hspace{27mm}
\mathfrak G^2\bigl(
Y_1\bigl(
\mathfrak G^3(
\mathfrak m_2(x_4,x_5))
\bigr)
\bigr)
\bigr)
\Bigr).
\end{aligned}
\]

明确这一点后，
考虑算子
$\mathfrak g(\mathbf Y)(\mathbf x)$。
按此记法，
算子在左，
$Y_i$ 在中，
$x_i$ 在右，
与式~\eqref{eq:7-4-79}
采用相同约定。
此外，
命题~\ref{prop:7-4-56}
中的符号，
就是在适当调整次数后，
改变算子、
$Y_i$ 与 $x_i$
的次序所得到的符号。
依照这些规则，
我们通过引理~\ref{lem:7-4-64}
按如下方式核对
命题~\ref{prop:7-4-56}
中的符号。

定义超交换子\liftIndex{I373}：
\[
[A,B]^{\mathrm{sup}}
=A\circ B
-(-1)^{\deg A\,\deg B}
B\circ A.
\]
